\chapter{Program, Scope, and Standing Conventions}
\label{comp:ch01}

\section{The question and the answer}
\label{comp:ch01-question}

The main manuscript asks two questions. What collective inference is available to a population of
agents, each a local section of a belief bundle and a model bundle of one principal bundle, all
placed at one base point $c_*$ and interacting through independent links under one normalized joint
law? And which parts of that inference survive when agents are coarse-grained exactly into
meta-agents? Its answer, announced in the opening of \MT{ch:introduction}, has four parts. The whole
population has one evidence bound, whose gap vanishes exactly when the recognition law equals the
posterior (\MT{thm:elbo-extended-gap}, \MT{cor:elbo-bound-tightness}). Each block of agents has a
local objective, a conditional free energy, such that when the block's conditional law is replaced
with the rest of the population held fixed, the change of the collective free energy equals the
outside-averaged change of the block's local objective (\MT{thm:obs-local-global-potential}); holding
the outside marginal fixed is load-bearing. Coarse channels preserve the evidence and split the free
energy exactly (\MT{thm:rg-exact-coarse-vfe}), with an exact effective action for what they retain
(\MT{thm:rg-effective-action}). And a finite-depth family of models forms successive meta-agents
under free-energy descent with learning, in a stated scope (\MT{thm:finite-depth-adaptive-formation}).
The program-aim section announces two further results on a homogeneous four-cycle: seeded formation
from every Gaussian initial law with source learning idle (\MT{prop:phil-seeded-contrast-formation}),
and law-level selection of one pairing by slow source learning for almost every Gaussian initial
mean (\MT{conj:phil-contrast-formation}, a theorem despite its label). The chapter also names what is
not done: how a parent's own dynamics are determined by the network it summarizes is
\MT{open:rg-pointwise-parent-dynamics}, which it calls the defining open problem of the present
phase.

\section{The picture}
\label{comp:ch01-picture}

Think of a lattice gauge theory on a graph whose sites all sit at one spacetime point. Each site
carries an agent. What the agent carries is not a field value but two probability laws, a belief
$q_i^b$ and a model $q_i^m$, on spaces on which the structure group $G$ acts through representations
$\rho_b$ and $\rho_m$; the agent also carries a frame $u_i$, a point of the one $G$-torsor over $c_*$
that says in which gauge the agent's laws are written. Each edge of the interaction graph carries a
group element $\Theta_e$ for each channel, an independent variable rather than something computed from
the frames, so products of links around loops (holonomies) can be nontrivial. Agents interact through
normalized record kernels $K_a$, which play the role of local factors of a Boltzmann weight.

From these local data the manuscript builds three global objects, in this order: one normalized joint
generative law $\mathbb P_{\theta,V}$ of everything, one recognition law $\mathbb Q_{V,o,X}$ of the
whole population, and one evidence bound $\Lelbo_V$. The order matters to the author, and it is the
first commitment of the program: every collective object is a complete probability object,
constructed from the agents' mechanisms and records or selected explicitly, never an extra agent added
to the population and never a sum of local objectives that happens to look like a global one.

\begin{analogy}{exact for the variational principle, structural for the gauge part}
The evidence identity below is the Gibbs variational principle of statistical mechanics: for any trial
state $q$, $\log Z=-F[q]+\KL(q\Vert p)$ with $F[q]=\E_q[\text{energy}]-\text{entropy}(q)$. Here the
``energy'' is $-\log$ of the joint generative density and $Z$ is the evidence. The gauge part is lattice
gauge theory over a zero-dimensional base: frames are local gauge choices, links are link variables,
holonomies are untraced Wilson loops. Two differences matter. The site variables are probability laws
on representation spaces, not field values, and the links need not carry a kinetic (plaquette) term;
whether they are fixed parameters or dynamical variables is declared model by model.
\end{analogy}

\section{Objects and notation}
\label{comp:ch01-objects}

\begin{decoder}
Agent $\mathcal A^i=(\mathcal C_i;q_i^b,q_i^m,u_i)$ & A belief section, a model section, and one common frame (\MT{def:geo-agent}); at one base point, two laws and a frame. & A site with its matter content and gauge choice. \\
Common frame $u_i$ & One frame for both channels; a gauge choice, not an assumption (\MT{prop:geo-common-frame-gauge}). & Gauge fixing at a site. \\
Cross morphism $\Phi$ & A map from belief fibers to model fibers (or back); extra structure, not implied by the bundles (\MT{def:geo-cross-morphisms}). & A coupling between two matter species. \\
Record kernel $K_a$ & A normalized conditional law on the scope of an interaction or observation. & A local factor of the Boltzmann weight. \\
Evaluator & The map that turns a model presentation $m_i$ into a generative kernel; a model is not itself a law. & The map from coupling constants to a Hamiltonian. \\
Constructed joint $\mathbb P_{\theta,V}$ & The normalized product of all kernels (\MT{prop:gen-exact-normalization}). & The Gibbs measure. \\
Selected coupling $\mathbb Q_{V,o,X}$ & One joint recognition law whose marginals are the agents' local laws, chosen by a declared measurable rule (\MT{hyp:prob-measurable-coupling-selector}). & A correlated trial state; not a product state unless chosen so. \\
Posterior slice $\boldsymbol\Pi_{\theta,V,o,X}$ & The exact posterior given the observation $o$ and the structural data $X$. & The exact conditioned Gibbs state. \\
Structural data $X$ & Graph, masks, and other structure, held as parameters unless a law for them is declared. & Quenched background. \\
\end{decoder}

A notation trap: the manuscript's $T_{ij}$ is a frame comparison, not a temperature or a transfer
matrix, and the formation chapters write their temperature as $\mathrm T$. The letter $\pi$ is both
the bundle projection and, with indices, the source prior of attention.

\section{The formulas, explained}
\label{comp:ch01-formulas}

The architecture display \MTeq{eq:intro-local-to-global-architecture} is a flow chart:
\[
\{\mathcal A^i,G_{\theta,i},\mathbb Q_{i,o,X},K_a\}_{i\in V,a\in\mathfrak A}
\longrightarrow
\bigl(\mathbb P_{\theta,V}\text{ by construction},\ \mathbb Q_{V,o,X}\in\operatorname{Cpl}(\mathbb Q_{\bullet,o,X})\text{ by selection}\bigr)
\longrightarrow\Lelbo_V .
\]
On the left are the local data: agents, their generative mechanisms, their local recognition laws,
and the record kernels. The joint law is built by multiplying normalized kernels; the population's
recognition law is chosen among the couplings $\operatorname{Cpl}$ of the local laws, the joint laws
with those marginals. The selection step is not trivial. Once two agents have non-degenerate laws, the
marginals do not determine the joint, and the correlations the selection puts in are exactly what a
mean-field (product) choice leaves out.

The one identity \MTeq{eq:elbo-intro-identity},
\[
\log p_{\theta,V}(o\given X)=\Lelbo_V(\mathbb Q_{V,o,X};X,o)+\KL\bigl(\mathbb Q_{V,o,X}\big\Vert\boldsymbol\Pi_{\theta,V,o,X}\bigr),
\]
reads as follows. The left side is the log evidence, the log partition function of the system
conditioned on the observed data. The first term on the right is the evidence lower bound, minus the
variational free energy of the population's recognition law. The second term is the relative entropy
from that recognition law to the exact posterior: nonnegative, zero exactly when the two agree as
measures. Two consequences are spelled out in the chapter. Getting every single-agent marginal right
does not close the gap when the posterior is correlated (\MT{cor:elbo-bound-tightness}); the missing
piece is the correlation structure. And the identity does not change when agents re-express their
laws in other frames, provided the kernels transform covariantly
(\MT{prop:gen-product-evidence-invariance}, \MT{cor:elbo-independent-frame-invariance}): the evidence,
like a partition function, is gauge invariant. A subtler point concerns restricted families: if the
posterior is excluded from the recognition family, every individual gap is positive, but the best gap
over the family may still be zero, approached without being attained (\MT{sec:restrict-principle}).

\section{The standing conventions, and which of them matter}
\label{comp:ch01-conventions}

\MT{sec:intro-conventions} fixes defaults for the whole manuscript, so that later statements mention
only departures. The reading guide lists them. Three conventions are physics rather than hygiene.

The first is that the interaction graph at $c_*$ is the support of the source prior of attention, a
structural parameter that is the complete graph unless masked. Gibbs attention weights are then
strictly positive on every allowed edge whose comparison cost is finite (the extended-cost rows of
\MT{prop:full-graph-gibbs-rows} may put zero weight on an infinite-cost edge), and all dynamic structure
is carried by attention weights and link values. Creating or deleting edges is not part of the model;
it is \MT{open:obs-edge-creation}. Loosely, this fixes the set of possible couplings the way a
canonical ensemble fixes the particle number: only coupling strengths and link values move, never
the graph. Induced coarse multigraphs are derived objects.

The second is that three transports must stay distinct. A frame comparison $T_{ij}$ relates two frames
in the same fiber: because all agents sit over the same point, their frames are points of one torsor
and differ by a unique group element, with no path involved (\MT{prop:geo-frame-comparisons}). A link
$\Theta_e$ is an independent variable on an edge (\MT{def:geo-graph-links}); setting every link equal to
the frame comparison gives frame-aligned, holonomy-free links, and general links carry holonomy that no
frame change removes. Parallel transport along a curve in the base is not used at all. Graph links are
therefore not base transport, and graph holonomy is neither the holonomy of a base connection nor a
topological property of the bundle.

The third is that ``coarse'' has distinct meanings that the manuscript keeps apart: a Markov channel,
which discards sample information while preserving normalization; a hard restriction, which
substitutes a lower-dimensional configuration into an energy and defines a new model; and parameter
aggregation inside a closed family. \MT{sec:cg-three-operations}, where the three are defined, calls
them Markov pushforward, energy precomposition, and Galerkin restriction (the compression of an
operator, which is how aggregation acts on a Gaussian precision). They can agree in special
realizations (Gaussian ones are worked out in Supplement~G), and a theorem about one does not
transfer to another.

\begin{hygiene}
Standard Borel spaces, measurability of laws and kernels, and the finite positive evidence of admitted
observations are hygiene; every model the program uses satisfies them. Absolute continuity of the
recognition law with respect to the posterior is deliberately not a default: it is clause (H3) of
\MT{hyp:elbo-evidence-domain}, and when it fails the bound is $-\infty$. This matters in practice,
because a recognition law supported on a lower-dimensional set (for example, one that identifies
agents exactly) is singular and costs an infinite free energy. The decision that structural data are
parameters unless a law is declared matters too: every statement about random graphs, random
partitions, or random memberships needs that law to be named.
\end{hygiene}

\section{What the manuscript declines to claim}
\label{comp:ch01-scope}

The scope paragraph of \MT{sec:intro-conventions} declines several claims for the whole manuscript. No
index of the theory is identified with physical time: not the inference step, the flow parameter, the
coarse-graining depth, or the Fisher duration. The chapter states what such an identification would
need: a named target clock, a calibration map with an uncertainty model, causal and compositional
consistency, and evidence against alternative parameterizations. The base is not spacetime and
carries no Lorentzian structure. Agents, laws, and meta-agents are mathematical objects, and calling a
node an agent is a naming decision. No canonical rule for choosing a recognition coupling, a
partition, or a membership is asserted. No renormalization semigroup with rescaling, fixed points, or
beta functions is constructed in the main manuscript; that apparatus is in Supplement~B
(\cref{comp:sb}). Identifying any limiting object with a physical law is posed as
\MT{claim:physical-law-identification}. Each of these is declined, not refuted.

\section{The program's commitments}
\label{comp:ch01-commitments}

\begin{scorecard}
\MT{sec:intro-program-aim} & Every collective object is a complete probability object built from the agents' mechanisms and records or selected explicitly. & \CHOICE & Methodological; it excludes mean-field sums that only look global. \\
\MT{sec:intro-program-aim} & Every statement about a parent holds for the normalized parent law that coarse-graining produces. & \CHOICE & Rules out ``effective'' parents defined by hand. \\
\MT{sec:intro-program-aim} & Formation of a parent is a theorem about a declared family, with its scope stated. & \CHOICE & The formation results are held to this; see \cref{comp:ch13,comp:ch14}. \\
\MT{sec:intro-conventions} & One base point; no base maps, neighborhoods, or connection descent. & \CHOICE & A phase decision; base transport is deferred to Supplement~B. \\
\MT{sec:intro-conventions} & Interaction graph is the support of the source prior, complete unless masked; no edge creation. & \CHOICE & Edge creation is \MT{open:obs-edge-creation}. \\
\MT{sec:intro-conventions} & Frame comparison, graph link, and base transport are distinct. & \CHOICE & Links are independent variables; holonomy is network-loop holonomy. \\
\end{scorecard}

\section{Open ends}
\label{comp:ch01-open}

The chapter names one open problem as the defining one of the phase: how a parent's dynamics follow
from the network it summarizes (\MT{open:rg-pointwise-parent-dynamics}). The coarse-graining of
\cref{comp:ch09,comp:ch11} produces the parent's law at each instant exactly; what is missing is a
theorem that this law evolves as an autonomous process of its own. The other open ends of the program
are collected in \hyperref[comp:open-ends]{Open Ends of the Theory}.

\begin{verdictbox}{\PARTIAL}
The chapter claims four results and names one open problem. Three of the four are proved as stated:
the one evidence bound, the local-global potential identity, and exact coarse-graining with an
effective action. ``Exact'' there means an identity for one channel, not a closed renormalization
flow; the chapter does not claim a flow, and \cref{comp:ch11} explains what is and is not closed. The
fourth, formation, is proved in a narrower sense than the program aim suggests: the finite-depth
theorem relaxes a prepared state through a hierarchy of interaction strengths supplied by hand, and
the endogenous witness, seeded formation, and the pairing-selection theorem all live on the
four-cycle (\cref{comp:ch13,comp:ch14}); the selection theorem is the strongest learning result, a
law-level statement for almost every initial mean, and it amplifies an asymmetry of the initial law
rather than creating one. The program aim
that parents act back on their children through normalized laws is set up (\cref{comp:ch12}) but not
closed, because the parent's own dynamics is the open problem the chapter names. The chapter is
candid about this; its opening states the scope of the formation theorem and names the open problem.
The program-aim paragraph, which speaks of collectives forming through the agents' coupled inference,
describes the goal rather than the delivered result.
\end{verdictbox}
